\section{\heiti 引言}

\subsection{\heiti 研究背景与意义}

随着我国碳达峰、碳中和战略的推进，抽水蓄能电站作为电力系统重要的调节电源，装机规模持续扩大。国家能源局数据显示，截至2024年底，全国抽水蓄能装机容量已突破5,000万千瓦，"十四五"期间年均新增投产超600万千瓦\upcite{nea2024psh}。在项目投资决策阶段，财务评价模型是评估项目经济可行性的核心工具，直接决定数百亿级投资的经济效益判断。

然而，能源行业大量财务决策高度依赖Excel模型——经调研，某真实抽水蓄能电站财务评价模型包含14个工作表、57,577个非空单元格和42,109个公式，其中9,199次跨Sheet引用使人工追溯需要2-3人天，289个含错误公式（占比0.69%）在人工审计中极难全面检出。模型复杂度与可理解性之间的鸿沟，使得财务模型实质上成为"黑盒"：计算逻辑难以追溯，公式业务含义不透明，错误传播路径不可见。

公式黑盒问题的本质在于语法与语义的结构性解耦：同一SUM函数在不同行可能代表融资计算或折旧摊销，取决于所在行的业务上下文。纯公式语法分析无法携带业务语义，而人工理解又受限于跨Sheet引用的规模（9,199次），这一矛盾构成了本文研究的出发点。

将Excel财务模型从单点专家经验转变为可查询、可追溯、可审计的知识资产，是解决上述问题的关键。通过依赖图构建、语义分类和知识图谱技术，可将人工审计中的核心路径追溯压缩至秒级响应，为模型审计、错误根因定位和计算逻辑验证提供系统化支撑。

\subsection{\heiti 相关工作}

\textbf{Excel公式分析与程序理解。} 电子表格因其灵活性被广泛应用于金融、工程和科研领域，但其"平民编程"特性导致高错误率。Hermans等\upcite{hermans2014spreadsheet}对超过200个真实电子表格的分析表明，40\%的模型存在公式错误，且错误随规模增长呈指数级扩散。Jansen\upcite{jansen2004spreadsheet}提出基于控制流和数据流的程序理解方法，在公式重构和错误检测方面取得进展。然而，现有方法主要聚焦公式语法层面的正确性检查，缺乏对业务语义的自动识别能力——这正是本模型77.6\%的公式为SUM/IF/算术运算、公式语法本身不携带业务语义的核心挑战。

\textbf{知识图谱在垂直领域的应用。} 知识图谱技术已在设备运维、医疗、金融等领域证明其价值。Wang等\upcite{wang2025knowledge}提出基于动态知识图谱推理的变压器故障诊断方法，在实体关系推理上取得显著效果。Du等\upcite{du2024wind}专门针对风电设备构建了故障知识图谱，整合设备部件、故障模式等实体。然而，已有知识图谱研究面向设备运维和工业控制，未涉及财务模型领域——能源行业Excel财务模型作为核心决策载体，其结构化解析和知识化潜力尚未被充分挖掘。

\textbf{大语言模型与结构化数据查询。} Text2SQL等技术实现了自然语言到结构化查询的映射\upcite{spider2018}，但面向Excel财务模型的查询需理解计算逻辑和依赖关系，超出现有方法的处理范围。LLM在公式理解方面展现出一定潜力，但推理延迟和领域专业术语处理能力在实际审计场景中尚难满足实时性要求\upcite{huang2024large}。

综上，现有工作存在以下研究空白：（1）能源行业大量复杂Excel财务模型（典型规模：14个工作表、42,109个公式）作为核心知识载体，其结构化解析和知识化潜力尚未被充分挖掘；（2）现有方法缺乏对跨Sheet复杂引用关系的自动化追踪能力——本研究模型中存在9,199次跨Sheet引用，人工审计需要2-3人天；（3）Excel公式与业务语义之间的自动映射机制尚未建立——本研究消融实验表明，纯公式结构方法准确率仅20.7\%，说明语法不携带业务语义；（4）在能源财务模型场景下，纯公式结构方法与上下文规则方法的性能差异缺乏系统实证分析，"跨业务Sheet"现象对自动标注的影响尚无记录。本文正是针对上述空白展开研究。

\subsection{\heiti 本文工作与贡献}

本文提出基于知识图谱的Excel财务模型解析与语义追溯方法，以真实抽水蓄能电站财务评价模型为研究对象，主要贡献如下：

\begin{enumerate}
	\item \textbf{基于真实数据的深度模型分析：} 以包含42,109个公式的真实抽水蓄能财务模型为研究对象，系统揭示了资本金/全投资双视角结构和跨业务Sheet现象等此前未被报道的业务复杂性，为同类研究提供基准参考。
	\item \textbf{Excel解析与DAG构建方法：} 提出基于有向无环图的公式依赖图构建方法，成功构建包含62,160个节点、78,472条边的依赖图，验证无循环引用，支持最长730步的计算路径追溯。
	\item \textbf{公式语义理解消融分析：} 通过299条分层采样数据集上的系统消融实验，揭示纯公式结构方法（准确率20.7\%）与上下文规则方法（88.3\%）之间的巨大差距，从实证角度证明"公式黑盒"问题的结构性本质。
	\item \textbf{能源财务知识图谱与智能问答：} 构建包含5类实体、6类关系的知识图谱（42,138实体、65,333关系），支持财务指标计算路径的秒级追溯，问答意图识别综合准确率92.3\%，IRR路径追溯效率较人工提升超5,000倍。
\end{enumerate}

本文其余部分安排如下：第2节分析抽水蓄能财务模型的结构特征；第3节描述Excel解析与语义理解方法；第4节阐述知识图谱构建过程；第5节介绍智能问答系统；第6节报告实验与评估；第7节总结全文。
